\subsection{Монотонные функции. Существование односторонних пределов у монотонных функций.}

\subsubsection*{Монотонные функции}
Пусть $f(x)$ определена на $<a, b> \subset \mathbb{R}$.

\begin{tcolorbox}[title=Определения монотонности]
    \begin{itemize}
        \item $f(x)$ \textbf{возрастает (нестрого)} на $<a, b>$, если:
        \[ \forall x_1, x_2 \in <a, b>:\; x_1 < x_2 \Rightarrow f(x_1) \leq f(x_2) \quad (f \uparrow) \]
        
        \item $f(x)$ \textbf{возрастает (строго)} на $<a, b>$, если:
        \[ \forall x_1, x_2 \in <a, b>:\; x_1 < x_2 \Rightarrow f(x_1) < f(x_2) \]
        
        \item $f(x)$ \textbf{убывает (нестрого)} на $<a, b>$, если:
        \[ \forall x_1, x_2 \in <a, b>:\; x_1 < x_2 \Rightarrow f(x_1) \geq f(x_2) \quad (f \downarrow) \]
        
        \item $f(x)$ \textbf{убывает (строго)} на $<a, b>$, если:
        \[ \forall x_1, x_2 \in <a, b>:\; x_1 < x_2 \Rightarrow f(x_1) > f(x_2) \]
    \end{itemize}
\end{tcolorbox}

\textbf{Теорема.} Если $f(x) \uparrow$ на $<a, b>$, то:
\begin{enumerate}
	\item $\exists \lim_{x\to a+0} f(x)$ (конечный или бесконечный), причем $\lim_{x\to a+0} f(x) = \inf_{<a,b>} f(x)$.
	\item $\exists \lim_{x\to b-0} f(x)$ (конечный или бесконечный), причем $\lim_{x\to b-0} f(x) = \sup_{<a,b>} f(x)$.
\end{enumerate}

\begin{tcolorbox}[title=Доказательство (для предела справа в точке $a$)]
	Рассмотрим два случая поведения функции:
	\begin{enumerate}
		\item \textbf{Случай 1:} Функция $f(x)$ ограничена снизу на интервале $<a, b>$.
		Тогда по аксиоме полноты существует конечная точная нижняя грань:
		\[ A = \inf_{<a,b>} f(x) \in \mathbb{R} \]
		По определению инфимума:
		\begin{itemize}
		    \item $\forall x \in <a,b>: f(x) \ge A$.
		    \item $\forall \varepsilon > 0$, число $A + \varepsilon$ не является нижней гранью, значит $\exists x' \in <a,b>: f(x') < A + \varepsilon$.
		\end{itemize}
		В силу возрастания функции, для всех $x$, лежащих левее $x'$ (т.е. $a < x < x'$):
		\[ A \leq f(x) \leq f(x') < A + \varepsilon \]
		То есть $|f(x) - A| < \varepsilon$.
		Это и означает (если взять $\delta = x' - a$), что $\lim_{x\to a +0} f(x) = A$.

		\item \textbf{Случай 2:} Функция $f(x)$ не ограничена снизу на $<a,b>$.
		\[ \forall E > 0\; \exists x' \in <a,b>: f(x') < -E \]
		В силу возрастания, для всех $x \in <a, x']$:
		\[ f(x) \leq f(x') < -E \]
		Это означает, что $\lim_{x\to a+0} f(x) = -\infty$.
	\end{enumerate}
\end{tcolorbox}

\begin{tcolorbox}[title=Доказательство (для предела слева в точке $b$)]
	\begin{enumerate}
		\item \textbf{Случай 1:} Функция $f(x)$ ограничена сверху на интервале $<a, b>$.
		\[ B = \sup_{<a,b>} f(x) \in \mathbb{R} \]
		По определению супремума:
		\begin{itemize}
		    \item $\forall x \in <a,b>: f(x) \le B$.
		    \item $\forall \varepsilon > 0$, число $B - \varepsilon$ не является верхней гранью $\Rightarrow \exists x' \in <a,b>: f(x') > B - \varepsilon$.
		\end{itemize}
		В силу возрастания функции, для всех $x \in <x', b>$ (т.е. $x' < x < b$):
		\[ B - \varepsilon < f(x') \leq f(x) \leq B < B + \varepsilon \]
		$\Rightarrow |f(x) - B| < \varepsilon$. Это означает, что $\lim_{x\to b-0} f(x) = B$.

		\item \textbf{Случай 2:} Функция $f(x)$ не ограничена сверху на $<a,b>$.
		\[ \forall E > 0\; \exists x' \in <a,b>: f(x') > E \]
		В силу возрастания, для всех $x \in <x', b>$:
		\[ f(x) \geq f(x') > E \]
		Это означает, что $\lim_{x\to b-0} f(x) = +\infty$.
	\end{enumerate}
\end{tcolorbox}


\textbf{Обозначения:}
\begin{itemize}
    \item $B<a,b>$ — класс функций, ограниченных на $<a, b>$.
    \item $C<a,b>$ — класс функций, непрерывных на $<a, b>$.
\end{itemize}

\textbf{Следствия:}
\begin{enumerate}
	\item Аналогичные утверждения верны для $f(x) \downarrow$ на $<a,b>$ (пределы равны $\sup$ слева и $\inf$ справа).
	\item Если $f(x)$ монотонна на $<a, b>$, то:
		\begin{enumerate}
			\item В любой внутренней точке $x_0 \in <a,b>$ существуют конечные односторонние пределы:
			\[ \lim_{x\to x_0-0} f(x) \le f(x_0) \le \lim_{x\to x_0+0} f(x) \]
			\item Точки разрыва могут быть только 1-го рода (скачки).
			\item Множество точек разрыва не более чем счётно.
		\end{enumerate}
\end{enumerate}

\begin{tcolorbox}[title=Доказательство конечности пределов во внутренней точке]
	Пусть $x_0 \in <a,b>$ и $f(x) \uparrow$.
	Возьмем любые точки $x' \in (a, x_0)$ и $x'' \in (x_0, b)$.
	В силу монотонности для любого $x$ из окрестности $(x', x'')$ выполнено:
	\[ f(x') \leq f(x) \leq f(x'') \]
	Значит, функция ограничена в окрестности точки $x_0$.
	По доказанной теореме пределы слева и справа существуют и они конечны (так как ограничены значениями в соседних точках).
\end{tcolorbox}

\subsubsection*{Ограниченность функций. (хз почяему я написал это сюда, но пусть будет)}
\begin{tcolorbox}
	\textbf{Определение.} $f(x)$ ограничена сверху на $X$, если
	\[ \exists M: \forall x \in X, f(x) < M \]
\end{tcolorbox}
\begin{tcolorbox}
	\textbf{Определение.} $f(x)$ ограничена снизу на $X$, если
	\[ \exists m: \forall x \in X, f(x) > m \]
\end{tcolorbox}
\begin{tcolorbox}
    Если $f(x)$ принимает комплексные значения, то ограниченность $f(x)$ на $X$:
    \[ \exists M: \forall x \in X, |f(x)| < M \]
\end{tcolorbox}

\begin{tcolorbox}
	\textbf{Определение.} Для $f(x): X \to Y$:
	\[ \sup_X f(x) = \sup f(X) \]
\end{tcolorbox}
Число $M = \sup_X f(x)$, если:
\begin{enumerate}
	\item $ \forall x \in X\;\; f(x) \leq M $
	\item $ \forall M' < M\;\; \exists x \in X: f(x) > M' $
\end{enumerate}

\begin{tcolorbox}
	\textbf{Определение.} Если $f(x)$ не ограничена сверху на $X$, то $\exists \{ x_n \} \subset X $ такая, что
	\[ f(x_n) \xrightarrow[n\to\infty]{} +\infty \]
\end{tcolorbox}

\begin{tcolorbox}
	\textbf{Определение.} $f(x)$ принимает max в точке $x_0 \in X$, если
	\[ \forall x \in X\;\; f(x) \leq f(x_0) \]
	\[ f(x_0) = \max_X f(x) \]
\end{tcolorbox}

\subsubsection*{Односторонние пределы}

\begin{tcolorbox}
	\textbf{Определение. Предел справа (по Коши).} \\
	Пусть $f(x)$ определена на $(a, b)$.
	Число $A$ называется пределом $f(x)$ при $x \to a + 0$ (или $x \to a, x > a$), если:
	\[ \forall \varepsilon > 0\; \exists \delta > 0:\; \forall x \in (a, a+ \delta) \implies |f(x) - A| < \varepsilon \]
\end{tcolorbox}

\begin{tcolorbox}
	\textbf{Определение. Предел слева (по Коши).}
	\[ B = \lim_{x \to b-0} f(x) \Leftrightarrow \forall \varepsilon > 0\; \exists \delta > 0:\; \forall x \in (b - \delta, b) \implies |f(x) - B| < \varepsilon \]
\end{tcolorbox}

\begin{tcolorbox}
	\textbf{Определение по Гейне (для предела справа).} 
	\[ \exists \lim_{x \to a + 0} f(x) = A \Leftrightarrow \forall \{x_n\} \subset (a, b): x_n \to a \implies \lim_{n \to \infty} f(x_n) = A \]
\end{tcolorbox}

\begin{tcolorbox}
	\textbf{Определение по Гейне (для предела слева).} 
	\[ \exists \lim_{x \to b - 0} f(x) = B \Leftrightarrow \forall \{x_n\} \subset (a, b): (x_n \to b \land x_n < b) \implies \lim_{n \to \infty} f(x_n) = B \]
\end{tcolorbox}

\textbf{Утверждение.} Существование предела по Коши справа (или слева) равносильно существованию соответствующего одностороннего предела по Гейне.

\textbf{Теорема (Связь односторонних и двусторонних пределов).}
\[ \exists \lim_{x \to a} f(x) = A \Leftrightarrow
\begin{cases}
	\exists \lim_{x \to a + 0} f(x)\\
	\exists \lim_{x \to a - 0} f(x)\\
	\lim_{x \to a + 0} f(x) = \lim_{x \to a - 0} f(x) = A\\
\end{cases} 
\]

\begin{tcolorbox}[title=Доказательство]
	$\Rightarrow$: Пусть $A = \lim_{x \to a} f(x)$. 
	По определению: $\forall \varepsilon > 0\; \exists \delta_\varepsilon > 0: \forall x$
	\[ 0 < |x - a| < \delta_\varepsilon \Leftrightarrow x \in <a - \delta_\varepsilon, a> \cup <a, a + \delta_\varepsilon> \implies |f(x) - A| < \varepsilon \]
	Рассмотрим сужения на интервалы:
	\begin{itemize}
	    \item Если $x \in <a - \delta_\varepsilon, a>$, то неравенство выполнено $\Rightarrow A = \lim_{x \to a - 0} f(x)$.
	    \item Если $x \in <a, a + \delta_\varepsilon>$, то неравенство выполнено $\Rightarrow A = \lim_{x \to a + 0} f(x)$.
	\end{itemize}

	$\Leftarrow$: Пусть $\lim_{x \to a + 0} f(x) = \lim_{x \to a - 0} f(x) = A$.
	\begin{itemize}
	    \item $A = \lim_{x \to a + 0} f(x) \Rightarrow \forall \varepsilon > 0\; \exists \delta_1 > 0\; \forall x \in <a, a + \delta_1>: |f(x) - A| < \varepsilon$
	    \item $A = \lim_{x \to a - 0} f(x) \Rightarrow \forall \varepsilon > 0\; \exists \delta_2 > 0\; \forall x \in <a - \delta_2, a>: |f(x) - A| < \varepsilon$
	\end{itemize}
	Возьмем $\delta = \min(\delta_1, \delta_2)$. Тогда для любого $x$ из проколотой окрестности:
	\[ 0 < |x - a| < \delta \implies x \in <a-\delta_2, a> \cup <a, a+\delta_1> \]
	В обоих случаях $|f(x) - A| < \varepsilon$.
	Следовательно, $\lim_{x \to a} f(x) = A$.
\end{tcolorbox}
